\documentclass[12pt, a4paper]{article}
\usepackage[indonesian]{babel}
\usepackage[margin=2.5cm]{geometry}
\usepackage{enumitem}
\usepackage{hyperref}

\begin{document}

\title{Analisis Kasus Dugaan Pelanggaran Etika Akademik}
\author{Teosofi Hidayah Agung - 5002221132}
\date{\today}
\maketitle

\section{Kronologi Kasus Pelanggaran}
Berdasarkan informasi yang terhimpun dari berbagai media dan klarifikasi institusi, kronologi kasus ini adalah sebagai berikut:
\begin{itemize}
  \item \textbf{Awal Mula Kecurigaan:} Kasus ini mencuat ke publik sekitar akhir Mei 2026 ketika seorang periset Indonesia yang sedang mengikuti konferensi ilmiah internasional (salah satunya di Kopenhagen, Denmark) merasa curiga terhadap sekelompok WNI (termasuk Rifaldy Fajar dan rekan-rekannya).
  \item \textbf{Pengajuan Abstrak Masif:} Kecurigaan muncul karena kelompok tersebut menyodorkan jumlah abstrak yang sangat banyak (dilaporkan hingga 19 abstrak) dalam waktu yang sangat singkat ke dalam acara konferensi tersebut. Jumlah dan kecepatan ini dinilai sangat tidak masuk akal untuk standar penelitian ilmiah asli.
  \item \textbf{Pencatutan Afiliasi Palsu:} Setelah ditelusuri, para terduga pelaku menggunakan afiliasi kampus di Indonesia (seperti UNY, ITB, dan Universitas Internasional Semen Indonesia/UISI) tanpa izin resmi. Beberapa bahkan menggunakan nama laboratorium atau unit studi yang sama sekali tidak eksis di kampus tersebut (misalnya, \textit{Health Management Laboratory}).
  \item \textbf{Motif Travel Grant:} Motif utama dari pengiriman abstrak massal dan fiktif ini diduga kuat adalah untuk mengincar \textit{travel grant} (dana bantuan perjalanan). Dengan abstrak yang diterima (seolah-olah sebagai peneliti), mereka bisa bepergian ke luar negeri (Eropa/Jepang) secara gratis.
  \item \textbf{Klarifikasi Institusi:} Kasus yang viral di platform X (Twitter) dan Instagram ini memicu berbagai kampus seperti UNY dan UISI merilis pernyataan resmi yang membantah keterlibatan kampus. Diketahui para pelaku berstatus sebagai alumni atau mahasiswa yang tidak aktif, dan bukan dosen atau periset resmi institusi tersebut.
\end{itemize}

\section{Jenis Pelanggaran (Fraud) dan Klasifikasinya}
Tindakan yang dilakukan dalam kasus ini mencakup beberapa jenis pelanggaran etika akademik yang sangat serius:
\begin{itemize}
  \item \textbf{Fabrikasi Data (Fabrication):} Membuat-buat data, metodologi, dan hasil penelitian dari ketiadaan. Makalah yang disubmit diduga kuat di-\textit{generate} menggunakan bantuan \textit{Artificial Intelligence} (AI) secara penuh tanpa ada eksperimen atau analisis data nyata yang dilakukan.
  \item \textbf{Falsifikasi (Falsification):} Memalsukan identitas dan afiliasi institusi. Mengklaim berasal dari kampus atau laboratorium tertentu untuk meningkatkan kredibilitas tulisan.
  \item \textbf{Kepengarangan Fiktif/Tidak Sah (Ghost/Gift Authorship):} Memasukkan nama-nama ke dalam daftar penulis yang tidak memiliki kontribusi ilmiah atau kepakaran di bidang terkait. Misalnya, memasukkan nama mahasiswa dari jurusan Manajemen ke dalam riset bioinformatika genetik yang sangat spesifik.
  \item \textbf{Klasifikasi Pelanggaran: Berat.} Kasus ini diklasifikasikan sebagai pelanggaran akademik \textbf{berat}. Fabrikasi data yang dilakukan dengan sengaja dan sistematis untuk menipu komite konferensi internasional demi meraup keuntungan finansial (\textit{travel grant}) adalah bentuk penipuan yang fatal. Hal ini tidak hanya merusak kredibilitas individu, tetapi juga mencoreng nama baik institusi pendidikan dan ekosistem peneliti Indonesia di mata global.
\end{itemize}

\section{Identifikasi Kecacatan pada Paper yang Dilampirkan}
Jika kita membedah teks abstrak dari poster presentasi tersebut, terdapat banyak \textit{red flags} atau kecacatan fatal:
\begin{itemize}
  \item \textbf{Kejanggalan Afiliasi dan Latar Belakang:} Sahnaz V Putri ditulis berafiliasi dengan Universitas Internasional Semen Indonesia (UISI). Pada kenyataannya, pihak kampus mengonfirmasi ia adalah mahasiswa jurusan Manajemen yang mana sama sekali tidak relevan dengan kajian modifikasi RNA sirkuler atau epigenomik spasial.
  \item \textbf{Penggunaan Email Mahasiswa oleh Alumni:} Rifaldy Fajar mencantumkan alamat email \texttt{rifaldy.fajar@student.uny.ac.id} untuk publikasi yang diterbitkan pada tahun 2025. Padahal, ia merupakan alumni UNY angkatan 2014 yang sudah lama lulus, sehingga penggunaan status \textit{student} (mahasiswa) pada tahun tersebut adalah bentuk penyesatan identitas.
  \item \textbf{Afiliasi Fiktif/Tidak Terverifikasi:} Rifaldy Fajar juga mencantumkan afiliasi ``The Integrated Mathematical, Computational, and Data Science for BioMedicine Research Community Foundation, Jakarta Barat'', yang diduga kuat merupakan institusi fiktif yang sengaja dibuat agar terdengar relevan dengan topik riset.
  \item \textbf{Klaim Output yang Terdengar ``Terlalu Sempurna'':} Abstrak mengklaim mengembangkan model komputasi \textit{CircEpiNet} dan menghasilkan angka statistik yang sangat presisi, seperti nilai AUROC 0.87 (95\% CI: 0.84-0.90) untuk prediksi fenotipe seluler dan \textit{effect size} = 0.71 untuk \textit{circAPP hyper-methylation}. Menghasilkan angka \textit{Confidence Interval} yang begitu rapi dalam abstrak tanpa adanya rekam jejak riset laboratorium adalah pola khas teks halusinasi yang dihasilkan oleh \textit{Large Language Models} (seperti ChatGPT) saat diminta menulis abstrak ilmiah.
\end{itemize}

\section{Kajian Pustaka dan Justifikasi Kepercayaan Hasil}
\begin{itemize}
  \item \textbf{Kajian Pustaka Singkat:} Secara konseptual keilmuan, topik ini sebenarnya merangkai berbagai metode mutakhir yang memang ada di dunia nyata. \textit{Circular RNAs} (circRNAs) memang sedang diteliti karena perannya sebagai pengatur ekspresi gen dan hubungannya dengan penyakit neurodegeneratif. Penggunaan \textit{Single-nucleus RNA sequencing} (snRNA-seq) dipadukan dengan \textit{spatial epigenomics} dan \textit{Deep Neural Network} (DNN) adalah masa depan dari bioinformatika untuk memetakan penyakit Alzheimer. Kata kunci yang digunakan dalam abstrak ini sangat relevan dengan tren riset medis terkini.
  \item \textbf{Justifikasi Hasil (Trusted / Tidak Trusted):} \textbf{SANGAT TIDAK TERPERCAYA (Untrusted)}. Walaupun \textit{buzzwords} dan konsep medis yang dituliskan adalah nyata secara teoritis, \textbf{isi dan hasil dari penelitian ini adalah palsu}. Para penulis memanfaatkan kemahiran AI generatif untuk merajut istilah-istilah medis kompleks menjadi sebuah paragraf yang secara gramatikal dan tata bahasa ilmiah terlihat sangat meyakinkan. Kesimpulan dan metrik statistik seperti deteksi 14 circRNAs hanyalah karangan algoritma, bukan hasil pemrosesan data sungguhan. Mengingat rekam jejak fabrikasi dari para pembuatnya demi kepentingan \textit{travel grant}, dokumen ini tidak memiliki nilai ilmiah (pseudosains) dan tidak boleh dijadikan referensi sitasi manapun.
\end{itemize}

\end{document}